\section{Strategi Optimasi Parameter Variasional}

\subsection{Parameter-Shift Rule (Gradien Kuantum Eksak)}
Perhitungan gradien terhadap nilai ekspektasi energi merupakan komponen yang sangat krusial dalam algoritma \textit{Variational Quantum Eigensolver} (VQE) guna memastikan pembaruan parameter sirkuit dilakukan secara efisien. Tanpa adanya informasi gradien yang akurat, proses optimasi akan terjebak dalam pencarian acak yang tidak efisien atau mengalami kegagalan konvergensi pada lanskap energi yang memiliki dimensi tinggi. Penggunaan metode \textit{finite difference} pada sistem kuantum nyata sering kali menghadapi kendala ketidakstabilan numerik yang disebabkan oleh adanya \textit{hardware noise} saat menggunakan pergeseran parameter yang sangat kecil. Sebagai solusi fundamental, \textit{Parameter-Shift Rule} (PSR) memungkinkan perhitungan turunan analitik yang eksak dengan memanfaatkan pergeseran parameter yang relatif besar, sehingga lebih tahan terhadap fluktuasi statistik pada perangkat keras kuantum. Figure 10 mengilustrasikan mekanisme pergeseran parameter sebesar $\pm \pi/2$ pada fungsi biaya yang memungkinkan estimasi gradien dilakukan tanpa bergantung pada aproksimasi linear yang rentan terhadap galat.

Secara formal, nilai ekspektasi energi $E(\theta)$ didefinisikan sebagai $\langle \psi(\theta) | H | \psi(\theta) \rangle$, di mana turunan analitiknya dapat diperoleh melalui selisih antara dua nilai ekspektasi yang digeser secara simetris. PSR tidak memerlukan penambahan \textit{qubit} tambahan atau arsitektur sirkuit pendukung yang kompleks, melainkan hanya memerlukan eksekusi sirkuit yang sama dengan nilai parameter yang berbeda. Keunggulan utama dari metode ini adalah kemampuannya dalam memberikan nilai gradien analitik yang eksak, yang secara matematis diturunkan dari sifat aljabar operator uniter yang membangkitkan rotasi pada sirkuit. Dengan menggunakan pergeseran sebesar $s = \pi/2$, komponen gradien terhadap parameter $\theta$ dirumuskan secara presisi sebagai berikut:
\begin{equation}
    \frac{\partial E}{\partial \theta} = \frac{1}{2} \left[ E\left(\theta + \frac{\pi}{2}\right) - E\left(\theta - \frac{\pi}{2}\right) \right]
\end{equation} (1)
Implementasi dari formulasi PSR ini memberikan landasan operasional yang kokoh bagi algoritma optimasi berbasis gradien untuk bekerja pada sistem kuantum era NISQ dengan tingkat presisi yang setara dengan perhitungan analitik pada komputer klasik.

\subsection{Fondasi Optimasi Klasik: Gradient Descent}
Prinsip utama dari algoritma \textit{Gradient Descent} dalam optimasi parameter sirkuit kuantum adalah melakukan pembaruan nilai parameter secara iteratif dengan melangkah berlawanan arah dengan vektor gradien fungsi biaya. Strategi ini didasarkan pada asumsi bahwa arah negatif dari gradien merupakan jalur penurunan tercepat menuju titik minimum lokal pada permukaan energi yang bersifat mulus. Parameter \textit{learning rate} memainkan peran vital dalam menentukan sensitivitas serta besarnya langkah optimasi yang diambil pada setiap iterasi guna menjamin stabilitas proses konvergensi. Figure 11 menyajikan visualisasi mengenai permukaan energi fungsi biaya beserta lintasan langkah penurunan gradien yang bergerak secara sistemik menuju titik energi terendah di dalam ruang parameter. Metode ini menjamin bahwa sistem akan selalu bergerak menuju kondisi ekuilibrium yang lebih stabil selama fungsi biaya memiliki sifat diferensiabilitas yang baik di seluruh domain pencarian.

Meskipun memiliki landasan teori yang kuat, implementasi \textit{Gradient Descent} pada masalah portofolio finansial dengan jumlah aset yang besar menghadapi tantangan serius dalam hal skalabilitas komputasi. Untuk setiap langkah pembaruan gradien pada sistem dengan $p$ parameter, algoritma ini membutuhkan setidaknya $2p$ pengukuran nilai ekspektasi secara terpisah berdasarkan formulasi PSR. Biaya pengukuran yang bersifat linear terhadap jumlah parameter ini menjadi beban yang sangat berat bagi perangkat keras kuantum saat ini, terutama ketika jumlah aset yang dioptimasi mencapai skala industri. Inskalabilitas ini memicu kebutuhan mendesak untuk beralih dari metode gradien deterministik menuju pendekatan optimasi stokastik yang lebih hemat dalam penggunaan sumber daya pengukuran. Oleh karena itu, transisi menuju algoritma yang mampu memberikan estimasi gradien yang akurat dengan jumlah sampel yang lebih sedikit menjadi kunci keberhasilan dalam implementasi VQE skala besar.

\subsection{Simultaneous Perturbation Stochastic Approximation (SPSA)}
Algoritma \textit{Simultaneous Perturbation Stochastic Approximation} (SPSA), yang awalnya diperkenalkan oleh James Spall, menawarkan inovasi radikal dalam menanggulangi masalah skalabilitas pada optimasi dimensi tinggi. Berbeda dengan metode tradisional, SPSA secara konsisten hanya membutuhkan dua pengukuran nilai ekspektasi per iterasi, terlepas dari seberapa banyak jumlah parameter yang sedang dioptimasi di dalam sirkuit. Efisiensi luar biasa ini dicapai melalui konsep perturbasi simultan, di mana seluruh elemen dalam vektor parameter "diguncang" secara bersamaan ke arah acak dalam satu langkah evaluasi tunggal. Pengurangan drastis dalam beban jumlah tembakan (\textit{shots}) ke perangkat keras kuantum menjadikan SPSA sebagai pilihan yang sangat pragmatis untuk menjalankan eksperimen optimasi pada sistem kuantum nyata. Figure 12 menunjukkan perbandingan efisiensi jumlah pengukuran antara metode \textit{Gradient Descent} konvensional dengan metode SPSA seiring dengan meningkatnya jumlah parameter optimasi.

Mekanisme estimasi gradien dalam SPSA mengandalkan penggunaan vektor perturbasi acak yang dipilih dari distribusi biner Rademacher dengan nilai elemen $\pm 1$. Selisih nilai energi yang diperoleh dari dua titik perturbasi tersebut digunakan secara kolektif untuk memperbarui seluruh vektor parameter melalui satu estimasi gradien stokastik yang efisien. Berdasarkan hukum bilangan besar, meskipun estimasi gradien pada satu iterasi tunggal mungkin bersifat kasar, rata-rata dari gradien SPSA akan berkonvergensi menuju gradien asli seiring dengan bertambahnya jumlah iterasi. Proses pembaruan parameter dilakukan mengikuti aturan iteratif $\theta_{k+1} = \theta_k - a_k \hat{g}_k$, di mana variabel kontrol $a_k$ berfungsi sebagai \textit{learning rate} dan $c_k$ sebagai besarnya langkah perturbasi. Kedua hiperparameter tersebut harus dirancang sedemikian rupa agar meluruh secara bertahap guna menjamin konvergensi yang stabil serta ketahanan yang tinggi terhadap \textit{noise} yang inheren pada data pasar finansial maupun output perangkat keras kuantum.
